\section{Introduction}
\label{sec:intro}

The pre-train-and-fine-tune paradigm has catalyzed the development of highly specialized neural networks, each trained to excel at a distinct task \cite{devlin2018bert, radford2019language, vit}. However, deploying an array of individual large-scale models in production leads to linear scaling of computational costs and memory overhead. To mitigate this, model merging has emerged as a compelling area of study, aiming to combine multiple task-specific expert weights into a single, multi-task network without requiring retraining or joint training from scratch \cite{wortsman2022model, task_arithmetic}. 

Early methods such as task arithmetic \cite{task_arithmetic}, Fisher merging \cite{matena2022merging}, and TIES-Merging \cite{yadav2023ties} rely on static, linear combinations of weights. While highly parameter-efficient and zero-cost during inference, static merging is inherently task-agnostic: once merged, the network uses the exact same weight configuration for all incoming inputs. To overcome this, dynamic model merging has been proposed to predict input-dependent routing coefficients $\alpha_k(x)$ at runtime, blending task-specific parameters dynamically on a per-sample or per-batch basis \cite{routing_soup}.

A prominent family of methods, such as the Layer-wise Low-dimensional Classical Router (\textbf{L3-Router}) \cite{l3_router_orig}, learns independent routing networks for each of the $L$ layers. In post-hoc dynamic ensembling, this unshared architecture is widely adopted but introduces two severe empirical flaws:
\begin{enumerate}
    \item \textbf{Layer-to-Layer Coefficient Ruggedness and Representation Drift:} When routing parameters are fully independent across sequential layers, the optimization degrees of freedom scale linearly with $L$. On small calibration splits, this can cause the coefficients to diverge abruptly layer-by-layer. This introduces high-frequency fluctuations in weight blending across sequential layers, causing cascading representation drift and distorting representation features as they propagate through the deep network.
    \item \textbf{Parameter Scaling Excess:} Learning independent routers for every layer (e.g., 240 parameters for a 12-layer backbone) overfits heavily to low-noise features and collapses catastrophically on noisy out-of-distribution (OOD) domains, such as the Street View House Numbers (SVHN) dataset, when calibration data is scarce (64 samples).
\end{enumerate}

To bypass these failure modes, alternative approaches like Quantum Wavefunction Superposition Merging (\textbf{QWS-Merge}) \cite{qws_merge_orig} have proposed non-monotonic, wave-inspired cosine activations, but we expose that QWS-Merge is overparameterized and has high complexity compared to our block-wise sharing alternatives.

True to the core tenets of \textbf{The Empiricist}, we reject speculative mathematical metaphors and present a foundational empirical deconstruction of dynamic routing mechanics. We evaluate our methods inside our controlled, high-throughput \textbf{Task-Conflict Model-Merging Sandbox}. To resolve sequential representation drift and parameter excess, we introduce the \textbf{Block-wise Weight-Sharing Router} (\textbf{BWS-Router}), which groups the $L$ layers into $G = L / M$ block groups and shares routing weights within each block, providing extreme parameter compression and inference-time computational savings.

Furthermore, we extend our evaluation to a **true physical sequential weight-space model-merging experiment** on PyTorch multi-layer MLP experts. In this physical setup, we propagate representations sequentially through layers whose weights are physically blended at runtime, eliminating any virtual-layer averaging and fully validating BWS-Router under realistic sequential deep propagation.

Through a large-scale grid search of over 1,280 experiment runs across 5 independent seeds, we systematically map the parameter-performance landscape. Our key contributions are:
\begin{itemize}
    \item We mathematically formalize and empirically prove the phenomenon of \textbf{coefficient ruggedness} in unshared dynamic model merging, showing that block-wise parameter sharing acts as an effective stabilizer against sequential representation drift.
    \item We introduce the \textbf{BWS-Router} architecture, which achieves an outstanding dynamic Joint Mean accuracy of \textbf{79.57 $\pm$ 1.14\%} across MNIST, FashionMNIST, CIFAR-10, and SVHN, using only 80 trainable parameters---a 66.7\% parameter reduction over unshared baselines, while the static average baseline completely collapses to \textbf{23.56 $\pm$ 2.91\%} under semantic conflicts.
    \item We validate BWS-Router under a **physical sequential weight-space model-merging framework**, demonstrating that our proposed shared router ($M=3$) achieves **45.26 $\pm$ 10.11\%** Joint Mean accuracy, completely outperforming static uniform physical merging (**17.88 $\pm$ 3.78\%**), and dramatically boosting heterogeneous mixed-batch accuracy by **+10.93\%** absolute over the unshared baseline (achieving **43.20 $\pm$ 22.49\%** vs. **32.27 $\pm$ 21.28\%**).
    \item We expose the extreme vulnerability of wave-based routing (QWS-Merge) to seed variations, establishing regularized classical projections as a far superior and more stable alternative.
    \item We establish a clear, intellectually honest selection rule for gating activation functions: while Softmax gating is empirically superior for closed-world, mutually exclusive classification due to its "implicit sum-to-one regularization" and "winner-take-all logit-amplification", independent Sigmoidal routing is preferred for open-world, decoupled deployment (e.g. handling OOD deactivation and non-exclusive multi-task ensembling).
\end{itemize}
